% ===================================================================
%  Generado por Erlen
%  Compílalo en Overleaf, o localmente con dos pasadas de pdflatex.
% ===================================================================
\documentclass[aspectratio=169,11pt]{beamer}
\usetheme{Boadilla}
% El papel verde grisáceo y la banda clara del título no vienen de fábrica: se fijan a mano, como en Papel.
\definecolor{salviaPapel}{HTML}{EBEFE9}
\definecolor{salviaBanda}{HTML}{DEE7DA}
\definecolor{salviaTinta}{HTML}{2B2A24}
\definecolor{salviaAcento}{HTML}{33654A}
\setbeamercolor{background canvas}{bg=salviaPapel}
\setbeamercolor{normal text}{fg=salviaTinta}
\setbeamercolor{structure}{fg=salviaAcento}
\setbeamercolor{frametitle}{bg=salviaBanda,fg=salviaTinta}
% El filete verde que cierra la banda pediría redefinir la plantilla del frametitle: en el PDF la banda va sin él.
% En pantalla la letra es Source Sans; para que el PDF coincida, añade \usepackage[default]{sourcesanspro}.
\usepackage[utf8]{inputenc}
\usepackage[T1]{fontenc}
% Español: se carga solo si babel-spanish está instalado (Overleaf lo trae).
\IfFileExists{spanish.ldf}{\usepackage[spanish,es-nodecimaldot]{babel}}{}
\usepackage{amsmath,amssymb}
\usepackage[version=4]{mhchem}
\usepackage{booktabs}
\definecolor{erlenacento}{HTML}{33654A}
\usepackage{tikz}
\usetikzlibrary{arrows.meta, calc}
\renewcommand{\figurename}{Figura}
\renewcommand{\tablename}{Tabla}
\usepackage{siunitx}
\sisetup{per-mode=symbol, range-units=single}
\setbeamertemplate{navigation symbols}{}
% Algunos temas (Metropolis entre ellos) insertan solos una diapositiva al
% empezar cada sección. Erlen ya emite la suya, así que se desactiva la
% automática: de otro modo cada sección saldría dos veces y la numeración
% del PDF dejaría de coincidir con la de la app.
\AtBeginSection{}
\AtBeginSubsection{}

\title[Diseño experimental]{Diseño experimental}
\subtitle{Ejemplo didáctico · datos ilustrativos}
\date{\today}

\begin{document}

% --- diapositiva 1 ---
\begin{frame}[plain]
  \transdissolve[duration=0.25]
  \titlepage
\end{frame}
\note{%
  {\scriptsize\textbf{Tiempo previsto: 0:30 min}}\par\medskip
  Ejemplo didáctico. Sustituye contenido y datos por los de tu trabajo. No atribuir estas cifras a un experimento.\par
}

% --- diapositiva 2 ---
\begin{frame}{La pregunta que guía el trabajo}
  \transdissolve[duration=0.25]
  {\large ¿Qué diseño permite atribuir una diferencia a la intervención?}
  \medskip

  {\small Caso de uso ilustrativo. Define aquí el sistema, la comparación y el alcance.}
\end{frame}
\note{%
  {\scriptsize\textbf{Tiempo previsto: 1:30 min}}\par\medskip
}

% --- diapositiva 3 ---
\begin{frame}{Una estrategia explícita}
  \transdissolve[duration=0.25]
  \begin{figure}
    \centering
    \definecolor{sa0f}{HTML}{CFDCE6}
    \definecolor{sa1f}{HTML}{9DB7CC}
    \definecolor{sa2f}{HTML}{6B93B2}
    \definecolor{sa3f}{HTML}{396E98}
    \definecolor{sat1843819360}{HTML}{15181C}
    \definecolor{sat1226267613}{HTML}{FFFFFF}
    % Con babel-spanish las palabras largas se parten solas y el texto respira mejor.
    \begin{tikzpicture}[x=1cm,y=1cm,font=\footnotesize]
      \draw[fill=sa0f, draw=none] (0.000,0.000) -- (2.226,0.000) -- (2.518,-1.254) -- (2.226,-2.508) -- (0.000,-2.508) -- cycle;
      \draw[fill=sa1f, draw=none] (2.644,0.000) -- (4.870,0.000) -- (5.162,-1.254) -- (4.870,-2.508) -- (2.644,-2.508) -- (2.936,-1.254) -- cycle;
      \draw[fill=sa2f, draw=none] (5.288,0.000) -- (7.514,0.000) -- (7.806,-1.254) -- (7.514,-2.508) -- (5.288,-2.508) -- (5.580,-1.254) -- cycle;
      \draw[fill=sa3f, draw=none] (7.932,0.000) -- (10.157,0.000) -- (10.450,-1.254) -- (10.157,-2.508) -- (7.932,-2.508) -- (8.224,-1.254) -- cycle;
      \node[text=sat1843819360, align=center, text width=1.75cm, inner sep=1pt] at (1.113,-1.254) {\bfseries Definir unidad experimental};
      \node[text=sat1843819360, align=center, text width=1.75cm, inner sep=1pt] at (3.976,-1.254) {\bfseries Asignar condiciones};
      \node[text=sat1843819360, align=center, text width=1.75cm, inner sep=1pt] at (6.620,-1.254) {\bfseries Medir con controles};
      \node[text=sat1226267613, align=center, text width=1.75cm, inner sep=1pt] at (9.264,-1.254) {\bfseries Analizar según el plan};
    \end{tikzpicture}
  \end{figure}
\end{frame}
\note{%
  {\scriptsize\textbf{Tiempo previsto: 1:30 min}}\par\medskip
  Explica las decisiones y los controles, no solo los nombres de las técnicas.\par
}

% --- diapositiva 4 ---
\begin{frame}{Mostrar la evidencia}
  \transdissolve[duration=0.25]
  \begin{table}
    \centering
    \caption{Diseño didáctico; no estudio ejecutado.}
    \begin{tabular}{cc}
      \toprule
      \textbf{Concepto} & \textbf{Plan de ejemplo} \\
      \midrule
      Unidad experimental & Una preparación independiente \\
      Factor & Condición A o B \\
      Control & Misma preparación sin intervención \\
      Repetición técnica & Lecturas de la misma preparación \\
      \bottomrule
    \end{tabular}
  \end{table}
\end{frame}
\note{%
  {\scriptsize\textbf{Tiempo previsto: 1:30 min}}\par\medskip
  Datos o modelos didácticos: reemplazar antes de presentar resultados propios.\par
}

% --- diapositiva 5 ---
\begin{frame}{Interpretación y límites}
  \transdissolve[duration=0.25]
  \begin{itemize}
    \item Las lecturas repetidas no equivalen a muestras independientes.
    \item Planear controles y exclusiones antes de ver resultados.
    \item Justificar el tamaño muestral con el objetivo del estudio.
  \end{itemize}
\end{frame}
\note{%
  {\scriptsize\textbf{Tiempo previsto: 1:30 min}}\par\medskip
}

% --- diapositiva 6 ---
\begin{frame}{Antes de compartir}
  \transdissolve[duration=0.25]
  \begin{itemize}
    \item Aleatorización y posible cegamiento
    \item Confusores y efectos de lote
    \item Plan de análisis y límites de inferencia
  \end{itemize}
  \medskip

  {\small Conserva el proyecto JSON y revisa la exportación final.}
\end{frame}
\note{%
  {\scriptsize\textbf{Tiempo previsto: 1:30 min}}\par\medskip
}

\end{document}